\chapter{Study System and Methods}
\label{ch:study-system}

\section{The Westhaven estuary}
\label{sec:estuary}

The Westhaven estuary is a shallow, micro-tidal inlet of about
\(38\) square kilometres on a sheltered stretch of coast. Its mean depth
is \(2.4\) metres, the tidal range rarely exceeds forty centimetres, and
salinity rises from about \(14\) at the river mouth to \(26\) at the
seaward sill. Eelgrass covers roughly a fifth of the sea floor, in
meadows that range from a few hundred square metres to more than a
square kilometre \parencite{aagaard2021remineralisation}.

\begin{figure}[tb]
  \centering
  \begin{tikzpicture}[x=1mm, y=1mm, font=\footnotesize]
    \fill[accent!5] (0,0) rectangle (104,58);
    \fill[black!10]
      (0,58) -- (0,36) .. controls (14,40) and (22,30) .. (34,34)
      .. controls (46,38) and (52,28) .. (64,30)
      .. controls (76,32) and (84,44) .. (104,42) -- (104,58) -- cycle;
    \fill[black!10]
      (0,0) -- (0,14) .. controls (18,10) and (30,20) .. (46,16)
      .. controls (60,12) and (72,18) .. (104,10) -- (104,0) -- cycle;
    \draw[black!35] (0,36) .. controls (14,40) and (22,30) .. (34,34)
      .. controls (46,38) and (52,28) .. (64,30)
      .. controls (76,32) and (84,44) .. (104,42);
    \draw[black!35] (0,14) .. controls (18,10) and (30,20) .. (46,16)
      .. controls (60,12) and (72,18) .. (104,10);
    \foreach \x/\y/\n in {8/30/1, 15/22/2, 22/28/3, 30/19/4, 38/30/5,
        46/21/6, 55/27/7, 62/19/8, 70/26/9, 78/20/10, 86/31/11, 94/23/12} {
      \fill[accent] (\x,\y) circle (0.9);
      \node[anchor=south, inner sep=0pt, yshift=1.5mm, black!75] at (\x,\y) {\tablefigures\n};
    }
    \node[black!60, anchor=west] at (2,52) {\itshape river mouth};
    \node[black!60, anchor=east] at (102,5) {\itshape seaward sill};
    \draw[black!60] (80,49) -- (100,49);
    \foreach \x in {80,90,100} \draw[black!60] (\x,48.2) -- (\x,49.8);
    \node[black!60, anchor=south] at (90,49.5) {2 km};
  \end{tikzpicture}
  \caption{Schematic plan of the Westhaven estuary with the twelve
    monitored meadows, numbered from the river mouth to the sill.}
  \label{fig:map}
\end{figure}

The twelve meadows shown in \cref{fig:map} were chosen to span the
salinity and depth gradients of the estuary. \Cref{tab:sites} lists their
main characteristics; meadows nearer the river mouth are shallower, more
turbid and, as \cref{ch:heatwave-response} shows, warmer in summer.

\begin{table}[tb]
  \centering\small
  \caption{Characteristics of the monitored meadows. Depth is relative to
    mean sea level; carbon stock refers to the upper 30 cm of sediment.}
  \label{tab:sites}
  \begin{tabular*}{\linewidth}{@{\extracolsep{\fill}}l r r r r@{}}
    \toprule
    Meadow & Depth (m) & Area (ha) & Salinity & C stock (Mg ha\(^{-1}\)) \\
    \midrule
    1--3   & 1.1 & 18.4 & 15.2 & 41.7 \\
    4--6   & 1.6 & 36.9 & 18.8 & 52.3 \\
    7--9   & 2.2 & 61.0 & 22.1 & 58.9 \\
    10--12 & 2.9 & 44.5 & 25.4 & 47.6 \\
    \midrule
    All    & 2.0 & 160.8 & 20.4 & 51.8 \\
    \bottomrule
  \end{tabular*}
\end{table}

\section{Measuring heat dose}
\label{sec:heat-dose}

Temperature loggers fixed ten centimetres above the canopy recorded water
temperature every fifteen minutes. Following \textcite{marsh2021heatwaves},
the heat received by a meadow in a summer is summarised as its
cumulative degree-heating-days,
\begin{equation}
  H = \sum_{d \in \mathcal{S}} \max\bigl(0,\; \bar T_d - T_{90}(d)\bigr),
  \label{eq:heat-dose}
\end{equation}
where \(\bar T_d\) is the daily mean temperature, \(T_{90}(d)\) the
ninetieth percentile of the 1990--2019 climatology for that calendar day,
and \(\mathcal{S}\) the period from June to September. Only days within a
heatwave, as defined in \cref{sec:heatwaves}, contribute to the sum.

\section{Shoot censuses and sediment cores}
\label{sec:censuses}

Shoot density was counted every month from April to October in ten
fixed quadrats of \(0.25\) square metres per meadow. Sediment cores of
fifty centimetres were taken each October with a piston corer, sliced at
two-centimetre intervals and analysed for organic carbon after
acidification. Ages were assigned by excess lead-210 using the protocol of
\textcite{tetley2020dating}; detailed procedures for both methods are
given in \cref{app:protocols}.

Warming experiments used the open-top chambers described by
\textcite{lindqvist2024chambers}, which raise canopy temperature by
\(2.1 \pm 0.3\)\,\textdegree C without shading the plants.
